% !TEX root = transcript.tex
% =====================================================================
% Rehearsal-script PDF for the pre-defense deck. Self-contained: unlike
% the slide decks in this directory, this file carries its own
% \documentclass and preamble rather than \input-ing preamble.tex --
% that file is a beamer preamble (aspectratio=169) built for slides and
% does not apply to a portrait working document. Palette, base size and
% hyphenation list are duplicated from preamble.tex on purpose, the same
% way thesis_paper keeps its own copy of shared material so a single
% file travels alone (see thesis_presentation/README.md, "Why two
% documents" / the thesis_paper note under "Figures").
%
% Source of truth is transcript.md in this directory; this file is a
% typeset rendering of it and carries no content of its own. Regenerate
% by hand after editing the Markdown -- there is no automated build step
% from .md to .tex.
%
% Build:
%   latexmk -pdf transcript.tex
% =====================================================================
\documentclass[12pt,a4paper]{article}

\usepackage[T1]{fontenc}
\usepackage[utf8]{inputenc}
\usepackage{lmodern}
\usepackage[margin=25mm,headheight=15pt]{geometry}
\usepackage{xcolor}
\usepackage{amssymb}
\usepackage{upquote}
\usepackage{booktabs}
\usepackage{array}
\usepackage{ragged2e}
\usepackage{enumitem}
\usepackage{framed}
\usepackage{needspace}
\usepackage{microtype}
\usepackage{xparse}
\usepackage{csquotes}
\usepackage{fancyhdr}
\usepackage[hidelinks,bookmarksopen=false,pdftitle={Pre-defense transcript},%
  pdfauthor={Md. Sakif Khan}]{hyperref}

% --- palette: identical RGB values to preamble.tex, duplicated rather
% than shared so this file does not reach outside its own directory. --
\definecolor{agrdark}{RGB}{28,54,78}
\definecolor{agraccent}{RGB}{176,58,46}
\definecolor{agrgrey}{RGB}{110,110,110}

% --- straight "quotes" typed in the source become real typographic
% quotes without hand-converting every pair; nothing in this document
% needs a literal straight ". -----------------------------------------
\MakeOuterQuote{"}
\setquotestyle{american}

% --- ragged table cells: same convention as preamble.tex -- a narrow
% column justifies worse than it sets ragged. --------------------------
\newcolumntype{L}[1]{>{\RaggedRight\arraybackslash}p{#1}}
\renewcommand{\arraystretch}{1.25}

\raggedbottom
\clubpenalty=10000
\widowpenalty=10000

% Proper nouns and acronyms are never hyphenated -- same list as
% preamble.tex.
\hyphenation{WebQSP CWQ ComplexWebQuestions GraphRAG Freebase LangGraph
  McNemar Cypher AGR}

\setlength{\parindent}{0pt}
\setlength{\parskip}{5pt}

% --- running header: shows which section the reader is in. Uses the
% real TeX \markright mechanism (not a plain \renewcommand) because
% headings here are hand-built rather than real \section commands, and
% a plain macro would report whatever heading TeX's page builder had
% most recently *read*, not what is actually on the shipped page --
% those differ by a page or so due to lookahead. \mark is a node in the
% page's own content list, so the output routine resolves it correctly.
\pagestyle{fancy}
\fancyhf{}
\renewcommand{\headrulewidth}{0.4pt}
\renewcommand{\footrulewidth}{0pt}
\fancyhead[L]{\small\color{agrgrey}Pre-defense transcript}
\fancyhead[R]{\small\itshape\color{agrgrey}\rightmark}
\fancyfoot[C]{\small\color{agrgrey}\thepage}
\fancypagestyle{plain}{\fancyhf{}%
  \fancyfoot[C]{\small\color{agrgrey}\thepage}%
  \renewcommand{\headrulewidth}{0pt}}

% One whole slide section -- heading through its last director's note
% -- kept off two pages: everything from \begin{slide} to \end{slide}
% is typeset into a single minipage, which the page breaker treats as
% one atomic box. It either fits in the remaining space on the current
% page or the entire box relocates to the next one; nothing inside it
% can ever be split. That is the same mechanism a table or a tall
% figure relies on to avoid straddling a page, and it needs no
% \needspace heuristic: an oversized box that will not fit anywhere
% simply overflows visibly (Overfull \vbox) rather than failing
% silently, and no section here is remotely close to a full page.
%
% "N --- Title" is in the deck's ink, the timing at the margin, an
% underrule, a PDF bookmark. Starred form (\begin{slide}*) adds the
% accent star marking the three slides the timing table also bolds.
% The title and the (time) marker sit in fixed-width columns rather
% than sharing one \hfill-joined line: three of the twenty-three titles
% are long enough that a plain \hfill either collapses to no gap at all
% (title just barely fits the line) or wraps the time onto its own
% left-aligned line (title alone fills the line). Fixed widths make the
% time marker's position independent of how the title happens to break.
\newlength{\slidetimewidth}
\setlength{\slidetimewidth}{20mm}
\newlength{\slideheadgap}
\setlength{\slideheadgap}{4mm}
% \mark commands placed inside a minipage never reach the page builder
% -- a documented LaTeX limitation, not a lookahead problem this time --
% so \markright is issued after \end{minipage} in the *end* code, back
% in the main vertical list where marks work normally. \slidemark
% carries the heading text that far. Likewise minipage resets \parskip
% to 0 internally, which would silently collapse the blank line between
% paragraphs inside every speech/note box, so it is restored explicitly.
\NewDocumentEnvironment{slide}{s m m m}{%
  \par\vspace{9mm plus 2mm minus 2mm}%
  \def\slidemark{#2 --- #3}%
  \begin{minipage}{\linewidth}
  \setlength{\parskip}{5pt}%
  \phantomsection\addcontentsline{toc}{section}{#2 --- #3}%
  \noindent\begin{minipage}[t]{\dimexpr\linewidth-\slidetimewidth-\slideheadgap\relax}
    \raggedright{\color{agrdark}\Large\bfseries #2 --- #3}%
    \IfBooleanT{#1}{\,\textcolor{agraccent}{\ensuremath{\bigstar}}}%
  \end{minipage}%
  \hspace{\slideheadgap}%
  \begin{minipage}[t]{\slidetimewidth}
    \raggedleft{\color{agrgrey}\itshape\large(#4)}
  \end{minipage}\par
  \vspace{0.8mm}{\color{agrdark!55}\hrule height 0.5pt}\vspace{3.5mm}%
}{%
  \end{minipage}\par
  \markright{\slidemark}%
}

% A plain (non-numbered) named section: Backup slides, Anticipated
% questions, Recovery notes, Delivery.
\NewDocumentCommand{\plainhead}{m}{%
  \needspace{5\baselineskip}%
  \phantomsection\addcontentsline{toc}{section}{#1}%
  \par\vspace{9mm plus 2mm minus 2mm}%
  \noindent{\color{agrdark}\Large\bfseries #1}\par
  \vspace{0.8mm}{\color{agrdark!55}\hrule height 0.5pt}\vspace{3.5mm}%
  \markright{#1}%
}

% The words to say aloud: a blockquote rendered as an agrdark rule down
% the left edge, via framed rather than a plain \fbox so the rule can
% still be colour-scoped per instance. Every use sits inside a `slide`
% minipage now, so it never actually gets the chance to break across a
% page -- the enclosing box relocates as a whole before that happens.
\newenvironment{speech}{%
  \def\FrameCommand{{\color{agrdark}\vrule width 2.2pt}\hspace{3.2mm}}%
  \MakeFramed{\advance\hsize-\width\FrameRestore}%
  \noindent\ignorespaces
}{%
  \unskip\endMakeFramed\par\vspace{1.5mm}%
}

% A director's note attached to one slide: same mechanism, agraccent
% rule, one size down -- so it reads as an aside next to the words
% rather than as more of the words.
\newenvironment{note}{%
  \def\FrameCommand{{\color{agraccent}\vrule width 1.6pt}\hspace{3.2mm}}%
  \MakeFramed{\advance\hsize-\width\FrameRestore}%
  \small\noindent\ignorespaces
}{%
  \unskip\endMakeFramed\par\vspace{1.5mm}%
}

\begin{document}
\thispagestyle{plain}
\markright{Pre-defense transcript}

\begin{center}
  {\color{agrdark}\Huge\bfseries Pre-defense transcript}\\[3mm]
  {\color{agrgrey}\large Md.\ Sakif Khan}
\end{center}
\vspace{7mm}

Rehearsal script for \textbf{\texttt{pre-defense-0421052099.pdf}} --- 23
pages: a title, twenty-one body slides, a closing slide.

The four backup slides are a \textbf{separate file},
\texttt{pre-defense-0421052099-backup.pdf}. You do not present them. Open it
alongside the main deck and jump to a slide when a question calls for one;
the table near the end of this file maps each to its question.

\textbf{Budget: 24~min~56~s of speaking against a 25-minute limit.} That
leaves four seconds, and every row below is set to exactly what its own
words take at 93~wpm --- so there is no slack on any individual slide
either. Both numbers are a constraint, not a cushion.

The table read 24:26 until its rows were checked against the words above
them, and fourteen of the rows demanded a faster rate than this script's
own --- the echo-attractor slide worst, at 125~wpm against 93. Every
round's additions had been costed against the total and never against the
row they landed in, so the table stayed internally consistent while
drifting away from the speech it describes. It is re-derived from the words
on every edit now, and every addition since has been paid for out of
another row rather than out of the total: the relation-agnostic bound on
slide 8 came out of slides 4, 7 and 8 itself, and slide 3 came out of
slides 2, 4 and 9.

\textbf{Four seconds is not a margin.} Everything that could move to the
answers has already moved there, and what is left is spoken because it has
to be. A real cushion has to come out of what is said: about 60 words buys
40 seconds. The recovery note below names the slides that can give it
up --- 4, 7 and 9 --- and that is a decision about those slides, not about
this one. Note that all three have now given once, and 4 and 9 twice.

Times below are \textit{cumulative at the end of that slide}. If you are
more than 40 seconds past a marker, use the recovery notes at the bottom.

\begin{center}
\small
\begin{tabular}{@{}cL{82mm}rr@{}}
  \toprule
  \textbf{\#} & \textbf{Slide} & \textbf{Slide time} & \textbf{Cumulative} \\
  \midrule
  1  & Title                                & 0:19 & 0:19  \\
  2  & The problem                          & 0:42 & 1:01  \\
  3  & Why the model does this              & 0:55 & 1:56  \\
  4  & Where existing approaches stop       & 0:41 & 2:37  \\
  5  & Research questions                   & 0:53 & 3:30  \\
  6  & AGR: an explicit state machine       & 1:21 & 4:51  \\
  7  & Constrained tools                    & 0:37 & 5:28  \\
  8  & The Structural Verification Layer    & 2:06 & 7:34  \\
  9  & One claim, three routes              & 0:37 & 8:11  \\
  10 & Environment and question sets        & 1:05 & 9:16  \\
  11 & Making the comparison fair           & 1:12 & 10:28 \\
  12 & \textbf{Main results}                & 1:19 & 11:47 \\
  13 & Accuracy against cost                & 0:58 & 12:45 \\
  14 & \textbf{RQ1: hop count}              & 1:15 & 14:00 \\
  15 & The caveat I want to raise myself    & 1:17 & 15:17 \\
  16 & RQ2: groundedness                    & 1:10 & 16:27 \\
  17 & RQ2: what verification contributes   & 1:51 & 18:18 \\
  18 & \textbf{RQ3: ablation}               & 1:02 & 19:20 \\
  19 & The result I did not expect          & 1:15 & 20:35 \\
  20 & The echo attractor                   & 1:41 & 22:16 \\
  21 & The benchmark was wrong 57 times     & 0:50 & 23:06 \\
  22 & Contributions and limitations        & 1:45 & 24:51 \\
  23 & Thank you                            & 0:05 & 24:56 \\
  \bottomrule
\end{tabular}
\end{center}

The three \textbf{bold} slides are the ones the committee will actually
interrogate. If you are running long, take time from 3, 6, and 8 --- never
from 11, 13, 14, 17 or 18. That is the same list the recovery notes
protect, and 14 is on it for a different reason: skipping it hands the
clipping issue to them.

% =====================================================================
\begin{slide}{1}{Title}{0:19}

\begin{speech}
Good morning. I'm Sakif Khan. This is my pre-defense on Agentic Graph
Reasoning --- knowledge graph navigation with verification before the
answer is emitted. My supervisor is Dr. Sadia Sharmin.
\end{speech}

\begin{note}
\textit{Don't read the title aloud. It's on the screen.}
\end{note}
\end{slide}

% =====================================================================
\begin{slide}{2}{The problem}{0:42}

\begin{speech}
Language models answer factual questions fluently. The difficulty is that
they answer just as fluently when they do not hold the fact.

That isn't a rhetorical claim --- it's measured here. In my own
no-retrieval control on WebQSP, the model asserted 661 entities. 179 of
them do not exist anywhere in the knowledge graph. That's 27.1 percent,
and every one was stated without a hedge.
\end{speech}

\begin{note}
\textbf{Beat after "27.1 percent."} It is the number that sets up the
whole talk, and the next slide opens on the question it raises.
\end{note}
\end{slide}

% =====================================================================
\begin{slide}{3}{Why the model does this}{0:55}

\begin{speech}
So why does this happen at all?

The part of the definition that matters is that this isn't a bug that
escaped testing. It comes out of how these models are built, so the
response has to be architectural.

Three origins. The training data. The model itself, where knowledge sits
spread across the weights with no index, so nothing separates a memorised
fact from a plausible continuation. And the prompt.

Only the third is ours; the other two need retraining. That is where this
thesis lives.
\end{speech}

\begin{note}
\textit{Name the three, don't read them. The slide carries the wording.}
\end{note}
\end{slide}

% =====================================================================
\begin{slide}{4}{Where existing approaches stop}{0:41}

\begin{speech}
Four families, and each stops somewhere specific --- the table walks left
to right. The first three stop before reasoning begins, or at a radius
fixed in advance. Agentic navigation does interleave retrieval and
reasoning, which is the right move.

But all four share one gap: whatever the final generation call produces is
what gets emitted. Nothing checks what the answer \textit{asserts}.
\end{speech}

\begin{note}
\textit{Walk the table left to right with the pointer. Don't read the
cells verbatim.}
\end{note}
\end{slide}

% =====================================================================
\begin{slide}{5}{Research questions}{0:53}

\begin{speech}
Three questions.

RQ1: does agentic navigation actually improve multi-hop accuracy --- and
does the advantage grow with the number of hops?

RQ2: given a system that already navigates the graph, what does a
claim-level check against the traversed triples contribute?

RQ3: which components earn their cost, in accuracy and in tokens?

One framing note. RQ2 asks what verification \textit{contributes}, not
whether it reduces hallucination. That's deliberate. The answer has
several parts, and a yes-or-no question would have hidden most of them.
\end{speech}
\end{slide}

% =====================================================================
\begin{slide}{6}{AGR: an explicit state machine}{1:21}

\begin{speech}
AGR is an explicit state machine --- not a prompt loop. Six nodes.

The planner decomposes the question into ordered sub-objectives. The
explorer scores candidate edges and expands the frontier. The evaluator
decides whether the sub-objective has been met. The backtracker undoes a
bad expansion and bans the edge that caused it, so it can't re-take the
same wrong turn. The verifier is the contribution; more on it shortly. The
answerer emits only what survived.

There are three cycles --- the three arrows going back to the Explorer:
continue, backtrack, retry. All three are bounded by explicit budgets
rather than by model behaviour. That gives a termination guarantee that
doesn't depend on the model cooperating: every cycle passes through a
router that checks a monotone counter.
\end{speech}

\begin{note}
\textbf{Do not say "exactly two cycles."} The diagram has three arrows
returning to the Explorer and the audience is looking at it while you
speak. The thesis caption says two and names two, but the figure source
calls the third one a cycle in its own comment. Naming them after the
edge labels --- continue, backtrack, retry --- means counting the arrows
confirms the sentence.
\end{note}

\begin{note}
\textbf{If asked about budgets, go to backup page 2.}
\end{note}
\end{slide}

% =====================================================================
\begin{slide}{7}{Constrained tools}{0:37}

\begin{speech}
The agent never writes a graph query. It gets four operations with fixed
signatures --- relations, neighbours, an adjacency check, and entity
linking.

This matters for more than safety. Every operation is deterministic and
logged with its arguments and result, so the traversal is a record --- and
that record is what the verification layer checks against.
\end{speech}
\end{slide}

% =====================================================================
\begin{slide}{8}{The Structural Verification Layer}{2:06}

\begin{speech}
This is the contribution.

Before anything is emitted, the draft answer is split into atomic claims.
Each claim is checked against the triples the agent actually traversed.
Claims that can't be grounded trigger targeted re-exploration, or are
dropped. A claim the traversal itself grounds keeps those triples and
carries them back with the answer, so the system hedges rather than
asserts.

Four things I want to be precise about. First, why \textit{structural}
--- the check is against the graph the agent walked, not the model's
opinion of its own output. A model grading itself is not an independent
check.

Second, \textit{structural} is a bound, not a boast. The first two routes
test adjacency, not the relation --- any edge between the endpoints
passes, either direction, so a mother claim survives on a child edge.
Only route three reads what the claim says.

Third, "its evidence" is narrower than it sounds --- one route of three
records it, and the log keeps the count, not the triples. Slide 17 has
both bounds.

Fourth, a claim can be true and still be the wrong answer. Structural
grounding cannot catch that, and I'll show you exactly that failure
shortly.
\end{speech}

\begin{note}
\textbf{The last sentence is a promise. Slide 20 pays it off.}
\end{note}
\end{slide}

% =====================================================================
\begin{slide}{9}{One claim, three routes}{0:37}

\begin{speech}
This is the path a single claim takes. Three routes to being called
supported, and the ordering is about cost: only the third spends a model
call. The first two are pure lookups against a record we already have ---
and neither of them reads the relation.

The "+ evidence" label sits on the first route only.
\end{speech}
\end{slide}

% =====================================================================
\begin{slide}{10}{The environment and the question sets}{1:05}

\begin{speech}
The environment is a Freebase-derived graph: 2.6 million entities, 8.3
million triples, seven thousand distinct relations. It imports in 36
seconds, which matters only because it means the whole thing is
reproducible on one machine.

Questions: 400 each from WebQSP and ComplexWebQuestions. The important
column is reachability --- for every question I verified whether the gold
answer is actually reachable in this graph. 97 percent on WebQSP, 99.2 on
CWQ.

That number is the \textit{ceiling} on every accuracy figure I'm about to
show. It means when a system misses, the miss belongs to the system and
not to the environment.
\end{speech}
\end{slide}

% =====================================================================
\begin{slide}{11}{Making the comparison fair}{1:12}

\begin{speech}
Five systems: a parametric control, vector RAG, static GraphRAG,
Think-on-Graph as the agentic comparison, and AGR.

All five run on one frozen backbone, at temperature zero, under the same
25-call budget, on the same questions, against the same graph.

That's what lets me attribute differences to architecture rather than to
model capacity or to spend. One thing is \textit{not} equal, and the
slide says so: Think-on-Graph prunes from a narrower candidate set. That
cuts against it, not for it --- a thinner set is cheaper.

These benchmarks predate current models and may be partly memorised.
Without a no-retrieval control I couldn't separate what retrieval
contributes from what the model already knew.
\end{speech}

\begin{note}
\textbf{Say the width line, don't skip it.} The slide claimed a
retrieval-budget control it does not have; the Think-on-Graph baseline
section names the widths as the one place a reader should look first for
a confound, and the conclusion ranks it limitation 5. The numbers are on
the slide and in the answer below --- you do not have to recite
40/20/300/200 here.
\end{note}
\end{slide}

% =====================================================================
\begin{slide}*{12}{Main results}{1:19}

\begin{speech}
Here is the comparison.

AGR reaches 0.755 Hits@1 and 0.642 F1 on WebQSP, and 0.522 and 0.469 on
ComplexWebQuestions. That's ahead of every baseline on both datasets.

Two things I'd draw your attention to beyond the top line.

First, vector RAG on ComplexWebQuestions --- 0.203, \textit{below} the
no-retrieval control at 0.307. One verbalised triple cannot contain a
chain, so single-shot retrieval is worse there than not retrieving at
all. GraphRAG is beside it at 0.205, but its one-hop radius confounds the
paradigm, so the claim rests on vector RAG.

Second, the cost columns. AGR spends 4,511 tokens and 6.2 calls against
Think-on-Graph's 3,615 and 12.8. More tokens, but half the calls. I'll
come back to why calls are the number that matters.
\end{speech}

\begin{note}
\textbf{Slow down here. This is the slide they read while you talk.}
\end{note}

\begin{note}
\textbf{Do not pool the two retrieval baselines.} The table puts 0.203
and 0.205 side by side and the pooled version of this point is the one
that gets walked back: results.tex calls GraphRAG the weaker evidence of
the two and says the claim rests on vector RAG. The paper retracted the
pooled claim in its own words. If they press on GraphRAG, the answer
below has the strata.
\end{note}
\end{slide}

% =====================================================================
\begin{slide}{13}{Accuracy against cost}{0:58}

\begin{speech}
Plotted, the frontier depends on which cost you charge, and the two costs
I record disagree.

On tokens --- the axis here --- Think-on-Graph is \textit{not} dominated.
It buys lower accuracy at a genuinely lower token price, and any honest
reading has to say so. The interior point on WebQSP is static GraphRAG,
which the parametric control beats on both axes at once.

On calls, which is what the budget actually meters, Think-on-Graph is
dominated on both datasets.

I'm showing both because stating it on one axis alone would overclaim.
\end{speech}
\end{slide}

% =====================================================================
\begin{slide}*{14}{RQ1: accuracy against hop count}{1:15}

\begin{speech}
This is the direct answer to RQ1, and it's a shape rather than a
difference of means.

Look at ComplexWebQuestions on the right. AGR goes 0.46, 0.55, 0.57 as
questions get harder --- one hop, two hops, three or more. It is the only
system on that dataset that ends above where it started.

Three of the other four decay monotonically. Think-on-Graph is the
exception to the monotonicity but not to the direction --- it falls and
partially recovers, still 0.08 below its own one-hop score.

I want to be careful about the left panel. On WebQSP the three-or-more
stratum has n equals 4. I don't draw a conclusion from four questions, and
the thesis doesn't either.
\end{speech}

\begin{note}
\textbf{Volunteering the n=4 weakness pre-empts the obvious attack.}
\end{note}
\end{slide}

% =====================================================================
\begin{slide}{15}{The caveat I want to raise myself}{1:17}

\begin{speech}
I want to raise this before you do.

If you split the questions by whether Think-on-Graph finished inside the
shared 25-call cap, then on the questions it \textit{finishes},
Think-on-Graph is ahead of AGR. 0.852 against 0.788 on WebQSP. 0.629
against 0.607 on CWQ.

The entire aggregate margin comes from the questions it cannot finish ---
where it drops to 0.197 and 0.188, because it is cut off mid-search. It
gets clipped on 29 percent of WebQSP and 44 percent of CWQ.

So the honest claim is not that AGR reasons better per step. It's that
AGR completes its reasoning inside a fixed budget, and Think-on-Graph
frequently does not. That's an efficiency claim, and it's the one I'm
making.
\end{speech}

\begin{note}
\textbf{Deliver this as a strength. It is the most defensible slide in
the deck.}
\end{note}
\end{slide}

% =====================================================================
\begin{slide}{16}{RQ2: does anything ungrounded get asserted?}{1:10}

\begin{speech}
RQ2. Does anything ungrounded get asserted?

Pooling both datasets: AGR asserts 1,709 entities and zero of them are
absent from the graph. The parametric control asserts 1,001 and 22.1
percent of them don't exist.

That looks like a headline result for the verification layer. It
isn't --- and this is the finding I think matters most.

Think-on-Graph also reaches 0.0 percent. It has no verification layer at
all.

So zero ungrounded assertion is a property of \textit{navigating a graph}
--- if you can only name entities you actually visited, you can't invent
one. It is not a property of my verification layer, and I don't claim it
as one.
\end{speech}
\end{slide}

% =====================================================================
\begin{slide}{17}{RQ2: so what does verification actually contribute?}{1:51}

\begin{speech}
Which leaves the real question.

What it does not do: removing it changes accuracy by an amount I cannot
detect. p equals 1.0 on both datasets. It does not earn its place on
Hits@1 or F1, and I report that as a negative result.

What it does do: it withholds what it cannot ground. Removing the layer
drops the hedge rate 23.2 to 20.2 percent on CWQ, 8.5 to 8.0 on
WebQSP --- direction only; I have not tested that column. It attaches
supporting triples to the claims traversal grounds, and pairs the answer
with that evidence at emission.

The layer's case rests on auditability, not accuracy --- and I'll be as
plain about how far that goes. Two bounds, both on the slide. One route
of three records evidence; \texttt{verify\_connection} and entailment
accept with nothing attached. And the logger writes the \textit{count} of
supporting triples and drops the list. You can confirm these answers came
from a system that tracked its evidence; you cannot open the record and
inspect it.
\end{speech}
\end{slide}

% =====================================================================
\begin{slide}*{18}{RQ3: what each component earns}{1:02}

\begin{speech}
RQ3, and this is a four-condition ablation with paired McNemar tests
against the full system on stratified halves.

Read the p-value columns. Backtracking: 0.727 and 1.0. Verification: 1.0
and 1.0. Learned scoring versus embedding-only: 0.664 and 0.481. For
three of the four components I detect no accuracy effect at this sample
size.

I want to be careful with that phrasing --- that is "no detectable effect
at n equals 200," not "confirmed no effect." These are underpowered for
small differences and I say so.

One row is different. Removing the planner, on WebQSP, p equals 0.006.
\end{speech}

\begin{note}
\textbf{Pause on the 0.006. Then turn the slide.}
\end{note}
\end{slide}

% =====================================================================
\begin{slide}{19}{The result I did not expect}{1:15}

\begin{speech}
Removing the planner \textit{improves} WebQSP. Plus 0.083 F1, at p equals
0.006, while cutting tokens 31 percent and calls from 6.2 to 4.0. Better,
cheaper, and faster, by deleting a component I built.

On ComplexWebQuestions it trends the other way --- minus 0.047 F1 at p
equals 0.088. Not significant, but the opposite sign, and that's the
useful part.

The explanation is that WebQSP is predominantly one hop. Decomposing a
one-hop question invents a second sub-objective that sends the frontier
away from the answer it already had.

So the conclusion is that decomposition should be gated on question
structure rather than applied unconditionally. That's a design conclusion
the measurement forced on me, not one it confirmed.
\end{speech}

\begin{note}
\textbf{Own this. A negative result you can explain is stronger than a
positive one you can't.}
\end{note}
\end{slide}

% =====================================================================
\begin{slide}{20}{Every failure, read: the echo attractor}{1:41}

\begin{speech}
I read all 259 remaining failures and labelled them. Top categories,
pooled.

The one I want to name is sixth, with 13 cases.

The characteristic error of a graph navigator is not invention. It's what
I call the \textbf{echo attractor} --- the system returns a real,
grounded entity that sits one hop from the correct answer. Ask for a
director and get the film. Ask for a capital and get the country.

Here's why it matters, and this is the promise I made on slide seven: any
grounding check \textit{passes} it. The entity is real, it was traversed,
the triple exists. It is true and it is wrong. Verification cannot catch
this by construction.

Different systems fall into it together, so no evaluation treating them
as independent can see it. Rescore whenever a majority agree --- a
natural thing to want --- and this becomes apparent correctness. That is
the contribution: the mechanism, not the count.
\end{speech}

\begin{note}
\textbf{Do not say "it appears across systems, so it is a property of the
task rather than of AGR."} That is defensive where the thesis is
substantive, and nothing on the slide blames AGR for it. The claim is
about evaluation: sec:echo calls the attractor "invisible to any
evaluation treating systems as independent", and section 1.6 says the
contribution is the mechanism "and what it means for consensus-based
evaluation, not the frequency". Slide 21 is the same finding seen from
the other side.
\end{note}

\begin{note}
\textbf{The table is pooled, and the slide says so.} Wrong and hedge are
never pooled in the thesis --- sec:taxonomy: "a pooled percentage would
describe neither" --- and pooling also hides the shape flip:
\texttt{composite\_claim} is 1 on WebQSP against 46 on CWQ. The caption
carries both facts and points at backup page 4, so the split census is
something you offer rather than something you are corrected with.
\end{note}

\begin{note}
\textbf{The full 12-category histogram is backup page 4 if anyone wants
it.}
\end{note}
\end{slide}

% =====================================================================
\begin{slide}{21}{The benchmark was wrong 57 times}{0:50}

\begin{speech}
Same pass, same cross-system agreement: consensus flagged 105 questions,
adjudication confirmed 41 --- and the gap is the attractor I just
described, not label noise.

Those 41 were excluded before the census. 17 more turned up inside it,
one is in both, so 57 distinct questions where the benchmark was wrong,
not the system.

Published defect rates for these benchmarks are rare and usually
anecdotal. This is a counted rate with a documented adjudication
procedure and per-question provenance.
\end{speech}
\end{slide}

% =====================================================================
\begin{slide}{22}{Contributions, and what I would not claim}{1:45}

\begin{speech}
To summarise. Six contributions, and they are the six the thesis claims
in section 1.6. The one I would underline is stratum-dependent
decomposition: the literature treats decomposition as straightforwardly
beneficial, and it is not.

The sixth is worded \textit{pre-registered} in the thesis; nothing was
filed with a registry, so I say \textit{pre-specified}.

And the limitations, which I'd rather state than be asked.

The first is the most serious. The verifier persists only what it
\textit{rejects}, so wrongful acceptance has no rate at all --- and the
same decision is why the output contract cannot be audited from the
record.

Then: no detectable accuracy gain from verification, and at n around 200
that is "not detected", not "none". Zero ungrounded assertion comes from
navigation, not my layer. One environment, one backbone, one annotator.
And Think-on-Graph leads where it finishes, and prunes from a narrower
candidate set --- 40 and 20 against my 300 and 200 --- so its unclipped
subset understates it.
\end{speech}
\end{slide}

% =====================================================================
\begin{slide}{23}{Thank you}{0:05}

\begin{speech}
Thank you. I'm happy to take questions.
\end{speech}
\end{slide}

% =====================================================================
\needspace{5\baselineskip}
\phantomsection\addcontentsline{toc}{section}{Backup slides}
\par\vspace{9mm plus 2mm minus 2mm}
\noindent{\color{agrdark}\Large\bfseries Backup slides ---
  \texttt{pre-defense-0421052099-backup.pdf}}\par
\vspace{0.8mm}{\color{agrdark!55}\hrule height 0.5pt}\vspace{3.5mm}
\markright{Backup slides}

A separate document. Open it before you start, minimised or on a second
screen. \textbf{Everything in this script refers to backup slides by the
page number the viewer shows}, which is what this table lists. The file
opens on a title page, so the first backup slide is page 2. An ordinal
would resolve one short of every one of them: counted that way, the
fourth backup slide is hedging rather than the census.

\begin{center}
\small
\begin{tabular}{@{}cL{50mm}L{78mm}@{}}
  \toprule
  \textbf{Page} & \textbf{Contents} & \textbf{Use when asked} \\
  \midrule
  2 & Budget configuration and enforcement sites &
    "How do you guarantee termination?" \\
  3 & Which budgets actually bind &
    "Is the 25-call cap fair to Think-on-Graph?" \\
  4 & Full 12-category failure histogram &
    "What were the other failure modes?" \\
  5 & Hedging rates, all five systems &
    "Doesn't it just refuse more often?" \\
  \bottomrule
\end{tabular}
\end{center}

\textbf{Page 3 is the important one.} It shows AGR never reaches the call
cap --- 0.0 percent on both datasets --- which is precisely what makes
the comparison against a clipped Think-on-Graph legitimate rather than an
artefact of a cap chosen to suit AGR.

% =====================================================================
\plainhead{Anticipated questions}

\textbf{"Isn't the 25-call cap arbitrary, and doesn't it favour AGR?"}
It's the cap Think-on-Graph's own paper operates under. And AGR never
reaches it --- zero percent on both datasets (backup page 3). If I raised
the cap, AGR's numbers would not move; Think-on-Graph's would. I say that
in the thesis rather than leaving it for someone to find.

\textbf{"Do the categories look the same on both datasets?"}
\textit{(Slide 20 is pooled.)}
No, and that is the more interesting answer. The census is reported split
in the thesis and never pooled, because wrong and hedge describe
different failure semantics and the proportions differ sharply by
dataset. The clearest case is \texttt{composite\_claim}: 1 on WebQSP
against 46 on CWQ. WebQSP questions mostly are not compound, so the
category barely exists there; on ComplexWebQuestions it is the largest
single failure mode. A pooled percentage would describe neither dataset.
The full split is backup page 4.

\textbf{"Doesn't GraphRAG show the same thing?"} \textit{(Slide 12 puts
0.203 and 0.205 side by side. Do not pool them.)}
No, and the thesis says so rather than letting the two numbers be read
together. GraphRAG retrieves a one-logical-hop neighbourhood, so its fall
on CWQ confounds static retrieval with a radius I chose in advance. The
evidence that it is the radius is in the strata: GraphRAG scores 0.44 on
WebQSP's two-hop questions, which are largely mediator paths a one-hop
expansion does reach, against 0.16 on the CWQ two-hop stratum, where the
chains are genuine compositions it cannot reach at all. A retriever
failing uniformly on depth would not show that gap. There is a second
bound on it: it takes at most 100 edges per topic entity with no ordering
imposed, and on 72.5 percent of questions at least one topic entity
exceeds that degree, so on those it answers from an arbitrary sample.
Vector RAG carries the paradigm claim, because one verbalised triple
cannot contain a chain at any radius.

\textbf{"Did both systems see the same candidate sets?"} \textit{(The
sharper form of the cap question. Slide 11 raises it deliberately ---
answer it, don't deflect.)}
No, and it is the one thing I do not hold constant. Think-on-Graph keeps
40 relations per entity and 20 neighbours per relation --- the pruning
widths of the algorithm as published --- where AGR keeps 300 and 200.
Measured over the committed tool logs, the 40-relation cut binds on 31.6
percent of the 1,651 entities it expanded and the 20-neighbour cut on
32.8 percent of its neighbour calls; AGR's relation cap binds once in
3,097 expansions and its neighbour cap on 3.3 percent of calls. So the
cut is real and it is asymmetric. What it cannot do is rescue the budget
argument: a narrower candidate set makes each step \textit{cheaper}, so
it cannot explain why Think-on-Graph runs out of calls. What it does mean
is that the residual gap on the questions it \textit{finishes} is
measured against a system searching a thinner pool, so that figure is a
lower bound on what it could resolve at equal width --- not an estimate
of it. Re-running it at 300 and 200 is the first item in my future work,
and it is limitation 5 in the conclusion.

\textbf{"Your verifier doesn't verify the relationship --- any edge
between the two entities passes."}
Correct, and I would rather name it than defend it: it is the layer's
principal acceptance risk, and section 6.8 lists it first among the
mechanisms of wrongful acceptance. Both structural routes ignore the
relation and the direction, so a claim that X is Y's mother is certified
by an edge recording that X is Y's child. It is a deliberate consequence
of one design choice. The claim's relation is free text and Freebase
predicates are not --- "X played for Y" is realised through a roster
mediator carrying no predicate named for playing --- so matching on it
would reject true claims wholesale. The division of labour is that the
structural routes buy recall for "some supporting edge exists" and the
third route, the entailment check, buys the semantic precision. What I
will not do is state the exposure at the smaller of the two scales
available. The population at risk is not the rare
\texttt{verify\_connection} route but every claim those two routes accept
between them, and the log does not separate them: on test, of 2,008
accepted claims, somewhere between 39 and all 2,008 were certified without
any test of the asserted relation. That is an interval two orders of
magnitude wide and I report it as one rather than choose a point inside
it. Narrowing it needs the logging change --- persist accepted claims with
the triples that matched them --- which is the same future-work item as
the previous answer, and it is what would let the relation be checked
where the mediator schema makes that possible.

\textbf{"If verification doesn't improve accuracy, why keep it?"}
Because accuracy was never the only claim. It converts silent error into
an explicit hedge, and it pairs the answer with the triples that ground
it at the point of emission. I report the null on accuracy rather than
hiding it --- and I report the bounds on the auditability too, which is
the next question.

\textbf{"Show me one supporting triple, then."} \textit{(Expect this.
Slide 16 invites it.)}
I can't, from the committed record, and that is a limitation rather than
an evasion. \texttt{RunLogger} writes \texttt{n\_supporting\_triples} ---
an integer --- and discards the list, so no committed artifact in this
work contains a single supporting triple; every statistic I quote about
them comes from that counter. What I can show is the traversal record
each answer was produced against. Two things follow and I state both: the
pairing is real inside the run and unavailable afterwards, and one
polarity of the verifier's error is therefore unmeasured. Persisting
accepted claims with their matching triples is one logging change, and it
is the first item in my future work. I did not make it late because it
would separate the code from results already frozen against it.

\textbf{"Why isn't the five-system comparison one of your
contributions?"}
Because the thesis does not count it as one, and slide 22 is the thesis's
list --- section 1.6, one subsection per item. The comparison and the
hop-count shape are results the contributions rest on; they get slides 11
and 13, which is where the weight belongs. An earlier version of that
slide promoted both to contributions and dropped the ablation, the
decomposition finding and the protocol to make room --- still saying
"six". The thesis had already been audited for exactly that mismatch,
where its conclusion counted four against section 1.6's six.

\textbf{"How many questions is that hedge difference, and is it
significant?"}
Six, on CWQ: 23.2 percent of 198 against 20.2. The sets nest --- there is
no question the ablated run hedged on that the full system asserted
on --- so those are exactly the six the ablated run answered and the
layer declined to. None of the six came back correct: all six were
assertions that would have been wrong, which is the mechanism, seen at
the only place the design isolates it.

On WebQSP it is one question, and there the ablated run got it right. So
across the 398 paired questions correctness moved twice, once each way. I
would rather say that than be shown it. No, it is not
significance-tested; the ablation's McNemar test is on correctness, not
on the hedge column, and I claim the direction and nothing more.
\textit{Do not} offer the no-retrieval contrast here. Its 12.2 percent is
a hedge rate, not an error rate --- backup slide 5 has it in a column
headed "WebQSP hedge \%" --- and AGR hedges \textit{less} than it does,
8.2 against 12.2, so that comparison argues the opposite of what it looks
like. No-retrieval's actual error rate is 170 wrong out of the 351
questions it asserts on.

\textbf{"Does every accepted claim get evidence attached?"}
No --- one route of three does. Traversed adjacency attaches every
traversed triple joining the pair; \texttt{verify\_connection} and the
entailment fallback accept with nothing attached. On the 80-question
development set, 13 answers carry no supporting triples: ten are hedges
that asserted nothing, one had every claim rejected, and two asserted a
single claim certified by the route that records no evidence. Median
across all 80 is 3, mean 4.1, maximum 16 --- read from the counter, since
the triples are exactly what the log does not keep.

\textbf{"Your planner result says your own design is wrong."}
It says the planner is wrong \textit{for one-hop questions}, and WebQSP
is mostly one-hop. On ComplexWebQuestions the effect reverses in sign.
The conclusion is that decomposition should be gated on question
structure --- which is a finding, and one I'd have missed without the
ablation.

\textbf{"n=4 in the three-hop WebQSP stratum is meaningless."}
Agreed, and I don't draw a conclusion from it. The hop-count claim rests
on ComplexWebQuestions, where the strata are 137, 211, and 49.

\textbf{"How do you know the gold answers are reachable?"}
Verified per question against the graph before scoring --- 97.0 percent
on WebQSP, 99.2 on CWQ. It's the ceiling on every Hits@1 I report, and
it's stated as such.

\textbf{"Where do the topic entities come from --- doesn't the system
have to find them first?"}
They are given by the datasets. \texttt{use\_gold\_entities} is on for
every run I report, so the question's annotated mentions go to
\texttt{search\_entity}, which resolves each one to a graph node through
the three-stage resolver. Mention \textit{detection} is assumed;
mention-to-node resolution is not --- that part the system does. It is
limitation 7 in the conclusion: the accuracies I report presume a linking
step a deployed system would have to perform, and I do not measure that
step. What it is not is a between-system confound. The three systems that
touch the graph --- AGR, Think-on-Graph and GraphRAG --- all seed from
the same annotated mentions, and neither the parametric control nor
Vector-RAG ever sees them.

\textbf{"Nine of your failures are one bug. Isn't the census measuring
your implementation?"}
In part, and I would rather say which part. Nine of the 38
\texttt{decomposition\_error} cases carry the subtype
\texttt{extraction\_bug}, and they share one mechanism: the evaluator
resolves every gold value, the drafted sentence names them verbatim, and
\texttt{answer\_entities} then collapses to the sentence's grammatical
subject. WebQTest-1215 drafts all six of Stephen Covey's professions and
scores \texttt{['Stephen Covey']}. The reasoning was complete and the
verifier certified it; what failed is the step that reads entities back
out of the draft. That is limitation 8, and the direction is the part
worth saying: it depresses my \textit{own} reported accuracy, so on that
question shape the headline numbers are a floor rather than an estimate.
The fix is an instruction to the claim decomposer, not an architecture
change, and it is the first of the three repairs the census earned. The
category split is backup page 4.

\textbf{"Is this reproducible?"}
Every number in the thesis is generated from frozen run records by a
script; none is transcribed. Same for the three data figures in this
deck --- they are pulled directly from the thesis's own generated
sources.

% =====================================================================
\plainhead{Recovery notes}

If you hit \textbf{11:47 (slide 13) more than 40 seconds late}, compress
as follows.

\begin{itemize}[leftmargin=*,itemsep=4pt,topsep=2pt]
  \item Slide 17 --- cut to: \textit{"It doesn't improve accuracy. It
    converts silent error into an explicit hedge and attaches evidence.
    The case is auditability."} Saves \(\sim\)35~s.
  \item Slide 21 --- cut to: \textit{"Reading every failure also found 57
    questions where the benchmark was wrong, with documented
    provenance."} Saves \(\sim\)25~s.
  \item Slide 3 --- name the three origins and stop; skip the model
    sentence. Saves \(\sim\)20~s. Slide 4 has already been compressed this
    way and cannot give again.
\end{itemize}

Never compress 12, 14, 15, 18, or 19. Slide 15 in particular: skipping it
means a committee member raises the clipping issue instead of you, and it
lands very differently that way.

% =====================================================================
\plainhead{Delivery}

\begin{itemize}[leftmargin=*,itemsep=4pt,topsep=2pt]
  \item \textbf{Say the number, then what it means.} Not the reverse.
    "Twenty-seven percent of asserted entities don't exist --- the model
    is confidently inventing."
  \item \textbf{Three slides are negative results} (16, 17, 18). Deliver
    them at normal pace, not apologetically. A student who reports a
    clean null is more credible than one who reports only wins.
  \item The word is \textbf{hedge}, not "refuse." A hedge is a
    calibrated non-assertion.
  \item If you lose your place, the takeaway bar at the bottom of the
    data slides is your prompt --- read it aloud and continue.
\end{itemize}

\end{document}
